% \iffalse meta-comment ---------------------------------------------
% Copyright (C) Shanghai Jiao Tong University
% The definition in this file is referred to the Visual Identity System
% from Shanghai Jiao Tong University (SJTU). 
% See https://vi.sjtu.edu.cn for more information.
% 
% SJTUG implements the design but doesn't hold the copyright.
% Any commercial usage in this file should be acknowledged by
% the administration of SJTU.
% More information about the license,
% see https://vi.sjtu.edu.cn/index.php/articles/bulletin/16. 
% ------------------------------------------------------------------- \fi
% \iffalse
%<*package>
\NeedsTeXFormat{LaTeX2e}
\ProvidesPackage{sjtuvi}[2021-08-21 Visual Identity System library for sjtubeamer 1.3]
%</package>
% \fi
% \CheckSum{0}
% \StopEventually{}
% \iffalse
%<*package>
% ------------------------------------------------------------------- \fi
%
% \subsection{SJTU VI Library}
%   This library provides the definition of basic elements in SJTU VI, including color, logo and some shapes. 
%
%  \begin{macro}{\DefineOption}
%   Define the beamer option on the corresponding package, key and value.
%    \begin{macrocode}
\def\DefineOption#1#2#3{
    % #1: package
    % #2: key
    % #3: value
    \DeclareOptionBeamer{#3}{\expandafter\def\csname sjtubeamer@#1@#2\endcsname{#3}}
}
%    \end{macrocode}
%  \end{macro}
%
%  \begin{macro}{\EqualOption}
%   To check if the option on package, key is equal to value.
%
%   Here, a dummy trick is used to pass the if condition.
%   Since LaTeX handles \verb"\if" differently.
%   \verb"\iftrue" will eliminate the nearest \verb"\else" and \verb"\fi" but remains other extra \verb"\fi" and throws errors.
%   To avoid this, if the macro is expanded after \verb"\if", T=T, the condition holds and finish the current pair. And continues to process the real defined macro. This solution is somehow to combat against the LaTeX compiler. In modern LaTeX 3, it is not so nasty to deal with neseted conditions. 
%   \verb"\iffalse" doesn't need to be considered.
%    \begin{macrocode}
\def\EqualOption{TT\fi\ifEqualOption} % dummy trick to pass the condition test
\def\ifEqualOption#1#2#3{
    % #1: package
    % #2: key
    % #3: value
    \newif\iftempif
    \edef\tempvar{\csname sjtubeamer@#1@#2\endcsname}
    \def\tempvalue{#3}
    \ifx\tempvar\tempvalue\tempiftrue\else\tempiffalse\fi
    \iftempif
}
%    \end{macrocode}
%  \end{macro}
%
% \subsubsection{Load TikZ}
%   Load TikZ package and its related library: \verb"pattern.meta" provides the interface to define a pattern; \verb"fadings" provides the method to create a fading mask; \verb"decoration.pathmorphing" provides the interface to user-define a decoration.
%    \begin{macrocode}
\RequirePackage{tikz}
\usetikzlibrary{patterns.meta}
\usetikzlibrary{fadings}
\usetikzlibrary{decorations.pathmorphing}
%    \end{macrocode}
%
% \subsubsection{Color Definition}
%   The following color is defined by SJTU VI.
%   \begin{description}
%       \item[primary] The primary color, which influences the color of title and the background of title page.
%       \item[secondary] The secondary color, which influences the color of subtitle.
%       \item[tertiary] The tertiary color, which provides the color for the blocks.
%   \end{description}
%    \begin{macrocode}
\definecolor{sjtuRedPrimary}{HTML}{9E1F36}         %engred
\definecolor{sjtuRedSecondary}{HTML}{F28101}       %orange
\definecolor{sjtuRedTertiary}{HTML}{FED201}        %yellow

\definecolor{sjtuBluePrimary}{HTML}{004098}        %problue 
\definecolor{sjtuBlueSecondary}{HTML}{298626}      %lightgreen
\definecolor{sjtuBlueTertiary}{HTML}{004D4B}       %lightgray
%    \end{macrocode}
%
% \subsubsection{Declare Logo}
% \begin{macro}{\definelogo}
%    Define a masked logo to make its different color variants.
%    The first argument assigns the file name and the second sets the cropping mode for the logo:
%    \begin{itemize}
%      \item[h] The horizontal mode. The mask will be cropped by a horizontal rectangle in the center.
%      \item[v] The vertical mode. The mask will be cropped by a vertical rectangle.
%      \item[f] The full mode. This is suitable for a square logo. 
%    \end{itemize}
%    You should define a logo \verb"\definelogo{mylogo}{f}" then use it in the contents like:
%    \verb"\mylogo{white}", where the second parameter is the color you want.
%    \begin{macrocode}
\def\definelogo#1#2{
  % #1: file name
  % #2: mode: v|h|f
  \def\tempmode{#2}
  \if v\tempmode % vertical
    \begin{tikzfadingfrompicture}[name=#1]
      \fill[black] (-1,-1) rectangle (1,1);
      \node {\includegraphics[height=1.5cm]{#1.pdf}};
    \end{tikzfadingfrompicture}
    \expandafter\gdef\csname#1\endcsname##1{
      \begin{tikzpicture}
        \begin{scope}
          \clip (3.80,0.70) rectangle (6.25,9.30);
          \fill [##1,path fading=#1] (0,0) rectangle (10,10);
        \end{scope}
      \end{tikzpicture}
    }
  \else
    \begin{tikzfadingfrompicture}[name=#1]
      \fill[black] (-1,-1) rectangle (1,1);
      \node {\includegraphics[width=1.5cm]{#1.pdf}};
    \end{tikzfadingfrompicture}
    \if h\tempmode % horizontal
      \expandafter\gdef\csname#1\endcsname##1{
        \begin{tikzpicture}
          \begin{scope}
            \clip (0.70,3.80) rectangle (9.30,6.25);
            \fill [##1,path fading=#1] (0,0) rectangle (10,10);
          \end{scope}
        \end{tikzpicture}
      }
    \else % full
      \expandafter\gdef\csname#1\endcsname##1{
        \begin{tikzpicture}
            \fill [##1,path fading=#1] (0,0) rectangle (10,10);
        \end{tikzpicture}
      }
    \fi
  \fi
}
%    \end{macrocode}
% \end{macro}
%    Define the required logo.
%    \begin{macrocode}
\definelogo{sjtubadge}{f}
\definelogo{cnlogo}{h}
\definelogo{enlogo}{h}
\definelogo{dlogo}{f}
\definelogo{vlogo}{v}
%    \end{macrocode}
%
% Specially designed for sjtubg.
%    \begin{macrocode}
\begin{tikzfadingfrompicture}[name=sjtubg]
  \fill[black] (-1,-1) rectangle (1,1);
  \node {\includegraphics[width=1.5cm]{sjtubg.pdf}};
\end{tikzfadingfrompicture}
\def\sjtubg#1{
  \begin{tikzpicture}
    \begin{scope}
        \clip (0.9,3) rectangle (9.1,7.5);
        \fill [opacity=0.1,#1,path fading=sjtubg] (0,0) rectangle (10,10);
    \end{scope}
  \end{tikzpicture}
}
%    \end{macrocode}
% \subsubsection{Shape Declarations}
% \begin{macro}{stamp}
%	  Declare stamp pattern to make a stamp array.
%
%	  The newest version of TikZ provides the interface to user-define a pattern. Obeying compatibility philosophy, use \verb"\pgfkeyvalueof" interface to get parameters in a standard way. The unit is first tested in a standalone file and previewed by TikZEdt.
%    \begin{macrocode}
\tikzdeclarepattern{
  name=stamp,
  parameters={
    \pgfkeysvalueof{/pgf/pattern keys/size},
    \pgfkeysvalueof{/pgf/pattern keys/xshift},
    \pgfkeysvalueof{/pgf/pattern keys/yshift},
  },
  defaults={
    size/.initial = 5pt,
    xshift/.initial = 0pt,
    yshift/.initial = 0pt,
  },
  bottom left={(
    -0.5*\pgfkeysvalueof{/pgf/pattern keys/size}
      +\pgfkeysvalueof{/pgf/pattern keys/xshift},
    -0.4*\pgfkeysvalueof{/pgf/pattern keys/size}
      +\pgfkeysvalueof{/pgf/pattern keys/yshift}
  )},
  top right={(
    0.5*\pgfkeysvalueof{/pgf/pattern keys/size}
      +\pgfkeysvalueof{/pgf/pattern keys/xshift},
    0.4*\pgfkeysvalueof{/pgf/pattern keys/size}
      +\pgfkeysvalueof{/pgf/pattern keys/yshift}
  )},
  tile size={(
    \pgfkeysvalueof{/pgf/pattern keys/size},
    0.8*\pgfkeysvalueof{/pgf/pattern keys/size}
  )},
  code={
    \def\s{\pgfkeysvalueof{/pgf/pattern keys/size}}%
    \tikzset{x=0.5*\s,y=0.2*\s}
    \fill[xshift=\pgfkeysvalueof{/pgf/pattern keys/xshift},
      yshift=\pgfkeysvalueof{/pgf/pattern keys/yshift}] 
      (-0.25*\s,0) 
      -- (-0.17*\s,0.06*\s) 
      -- (-0.17*\s,0.1*\s) 
      -- (0.17*\s,0.1*\s) 
      -- (0.17*\s,0.06*\s)
      -- (0.25*\s,0) 
      -- (0.17*\s,-0.06*\s) 
      -- (0.17*\s,-0.1*\s) 
      -- (-0.17*\s,-0.1*\s) 
      -- (-0.17*\s,-0.06*\s) -- cycle;
  }
}
%    \end{macrocode}
% \end{macro}
%
% \begin{macro}{\stamparray}
%	  Create the stamp array in the TikZ environment.
%
%	  Notice \TeX{} is not good at handling parameters. Always remember to store it into a temporary variable. Register \verb"\pgfmathresult" will store the result of \verb"\pgfmathparse".
%    \begin{macrocode}
\providecommand{\stamparray}[3]{
  %#1: pattern size
  %#2: starting point
  %#3: ending point
  \usebeamercolor{palette primary}
  \fill [pattern={stamp[size=#1]},
    pattern color=bg!50!fg] #2 rectangle #3;
  \def\s{#1}%
  \pgfmathparse{0.5*\s}\let\xs=\pgfmathresult%
  \pgfmathparse{-0.4*\s}\let\ys=\pgfmathresult%
  \fill [pattern={stamp[size=#1,xshift=\xs, yshift=\ys]},
    pattern color=bg!50!fg] #2 rectangle #3;
}
%    \end{macrocode}
% \end{macro}
%
% \begin{macro}{stampline}
%	  Declare a decoration to make a loop stampline.
%
%	  Notice that \verb"auto corner on length" is open to avoid spikes where the state hasn't meet \verb"final" yet.
%    \begin{macrocode}
\pgfdeclaredecoration{stampline}{initial}
{
  \state{initial}[
    width=\pgfdecorationsegmentlength,
    auto corner on length=\pgfdecorationsegmentlength]
  {
    \def\l{\pgfdecorationsegmentlength}%
    \pgfpathlineto{\pgfpoint{0.25*\l}{0pt}}
    \pgfpathlineto{\pgfpoint{0.33*\l}{0.06*\l}}
    \pgfpathlineto{\pgfpoint{0.33*\l}{0.1*\l}}
    \pgfpathlineto{\pgfpoint{0.67*\l}{0.1*\l}}
    \pgfpathlineto{\pgfpoint{0.67*\l}{0.06*\l}}
    \pgfpathlineto{\pgfpoint{0.75*\l}{0pt}}
    \pgfpathlineto{\pgfpoint{\l}{0pt}}
  }
  \state{final}
  {
    \pgfpathlineto{\pgfpointdecoratedpathlast}
  }
}
%    \end{macrocode}
% \end{macro}
%
% \iffalse
%</package>
% \fi
%
% \Finale
\endinput
